\documentclass[11pt,a4paper,UTF8,fontset=none]{ctexart}
\setCJKmainfont[
  Path=assets/fonts/,
  AutoFakeBold=2.2
]{NotoSansSC.ttf}
\setCJKsansfont[
  Path=assets/fonts/,
  AutoFakeBold=2.2
]{NotoSansSC.ttf}
\setCJKmonofont[
  Path=assets/fonts/
]{NotoSansSC.ttf}
\usepackage[a4paper,margin=2.25cm]{geometry}
\usepackage{booktabs,longtable,array,tabularx}
\usepackage{graphicx}
\usepackage{xcolor}
\usepackage{amsmath,amssymb}
\usepackage{listings}
\usepackage{hyperref}
\usepackage{fancyhdr}
\usepackage{microtype}
\usepackage{enumitem}
\usepackage{caption}

\definecolor{navy}{HTML}{16324F}
\definecolor{blue}{HTML}{2D6A9F}
\definecolor{lightblue}{HTML}{EAF2F8}
\definecolor{warning}{HTML}{9C2F2F}
\hypersetup{colorlinks=true,linkcolor=navy,urlcolor=blue,citecolor=navy}
\setlist{nosep,leftmargin=2em}
\setlength{\parindent}{2em}
\setlength{\parskip}{0.25em}
\renewcommand{\arraystretch}{1.15}
\pagestyle{fancy}
\fancyhf{}
\fancyhead[L]{对称密码的软件实现与优化}
\fancyhead[R]{课程项目}
\fancyfoot[C]{\thepage}
\setlength{\headheight}{14pt}
\lstset{
  basicstyle=\ttfamily\small,
  keywordstyle=\color{blue},
  commentstyle=\color{gray},
  frame=single,
  breaklines=true,
  columns=fullflexible,
  backgroundcolor=\color{lightblue},
  xleftmargin=1em,
  xrightmargin=1em
}

\input{tmp/generated_results.tex}

\title{\bfseries 对称密码算法的软件实现与优化\\[0.4em]
\large AES／SM4／GIFT／TWINE 与 CTR／GCM／XTS}
\author{课程项目实现与实验报告}
\date{2026 年 7 月}

\begin{document}
\maketitle

\begin{abstract}
本项目从统一的 C11 标量基线出发，实现 AES-128/192/256、SM4-128、
GIFT-64（128-bit key）和 TWINE-80/128，并系统比较 T-table、
shuffle/bitslice/fixslice 与新指令集。工作模式覆盖全部算法的 CTR，以及
AES/SM4 的 GCM、XTS；GCM 使用 ARM PMULL 加速 GHASH，XTS 同时支持
IEEE~1619 与 GB/T~17964-2021 的 tweak 乘法约定及 ciphertext stealing。
Apple M2 Pro 与 Intel Core i9-13900H 上分别保存 3 次预热、15 次正式测量
的全部样本。x86-64 原生验证执行 AES-NI、VAES、SSSE3、GFNI、PCLMUL 和
VPCLMUL，包含 \XGenericChecks 项通用检查、\XBackendChecks 项专用后端检查
与 \XOpenSSLChecks 项 OpenSSL 差分；VSM4 因目标 CPU 不支持而明确跳过。
\end{abstract}

\noindent\textbf{关键词：} 对称密码；AES；SM4；GIFT；TWINE；T-table；
NEON TBL；PMULL；GFNI；CTR；GCM；XTS

\tableofcontents
\newpage

\section{材料、目标与交付边界}

课程材料为 86 页《对称密码的软件实现与优化》。原始 PPT 的项目副本
SHA-256 为
\begin{center}\ttfamily\small
69b70bead7ddf5feb3f12d582c5770c43\\[-0.2em]
241110eda6f45f9670cd4f18fc903cb
\end{center}
该值与用户提供文件逐字节一致；逐页主题与代码/实验落点见
\texttt{materials/SLIDE\_MAP.md}。材料从微架构、寄存器、T-table 讲到
shuffle、复合域、AES/SM4 指令、GIFT/TWINE fixslicing 与 CTR/GCM/XTS，
因此项目也按“基础实现 $\rightarrow$ 表优化 $\rightarrow$ SIMD/新指令
$\rightarrow$ 模式优化”组织。

\begin{quote}\color{warning}
\textbf{安全边界：}本项目用于教学和可复现实验，不是认证密码库。尤其
T-table 会以秘密值索引缓存，不应成为真实系统的默认后端。
\end{quote}

\section{微架构基础与优化模型}

\subsection{吞吐率不是指令条数}

现代处理器可多发射并行执行，实际耗时由执行端口、依赖链、指令时延和
吞吐共同决定。因而同样的轮函数可以有三类改进：
\begin{enumerate}
  \item 缩短关键依赖链，例如把互不相关的分组交错执行；
  \item 减少内存层次开销，例如把状态保持在寄存器，或缩小表；
  \item 用一条 SIMD/密码指令替代多条标量指令。
\end{enumerate}
全展开可以删除循环控制，却会扩大指令缓存占用。本项目保留清晰的循环基线，
只在数据表示和后端处做可测量的变换，同时以对象尺寸 CSV 记录代码代价。

\subsection{端序与未定义行为}

所有外部字节串都通过 \texttt{sc\_load\_be32/64} 与
\texttt{sc\_store\_be32/64} 显式转换。2~KiB SM4 重叠表虽然利用偏移
读取旋转后的 32-bit 字，却不把非对齐地址强制转换成 \texttt{uint32\_t*}，
从而避免别名和对齐未定义行为。这一点由 ASan/UBSan 与
\texttt{-Werror} 共同约束。

\section{四种分组密码的基础实现}

\subsection{AES}

AES 按 FIPS~197 实现 128/192/256-bit 密钥扩展、SubBytes、ShiftRows、
MixColumns 及逆变换。参考路径以 16 字节状态表示；T-table 同时实现正向与
逆向四表。ARM 路径用 AESE/AESMC 与 AESD/AESIMC，x86 路径用 AES-NI，
多块时可进入 VAES。所有路径使用同一轮密钥并与 FIPS KAT 交叉比较。

\subsection{SM4}

SM4 使用 32 轮广义 Feistel 更新：
\[
  X_{i+4}=X_i\oplus L\!\left(\tau(X_{i+1}\oplus X_{i+2}
  \oplus X_{i+3}\oplus rk_i)\right),
\]
\[
  L(B)=B\oplus(B\lll2)\oplus(B\lll10)\oplus(B\lll18)\oplus(B\lll24).
\]
密钥扩展和数据轮均依据 GB/T~32907-2016。实现把输入、输出与轮密钥端序
明确分开，并用官方向量
\texttt{012345...3210 $\rightarrow$ 681edf...4246} 验证。

\subsection{GIFT-64 与 TWINE}

GIFT-64 采用四路 bitslice 状态，S 盒直接变为四个位平面上的布尔网络，
位置换由固定掩码网络完成；批量 API 每次处理四个独立分组。TWINE 以
16 个 nibble 表示状态，分别实现 80-bit 与 128-bit
密钥调度；KAT 密文为 \texttt{7c1f0f80b1df9c28} 与
\texttt{979ff9b379b5a9b8}。

\section{统一接口与运行时分派}

公共上下文 \texttt{sc\_ctx} 对外保持不透明，提供密钥初始化、单块/多块
加解密、块长/算法/后端查询和明确错误码。单块 API 支持原地操作。
后端可选 \texttt{ref}、三种 \texttt{ttable}、\texttt{shuffle}、
\texttt{aes-hw}、\texttt{gfni}、\texttt{sm4-hw} 与 \texttt{auto}。

\begin{lstlisting}[language=C,caption={统一接口的核心调用}]
sc_ctx ctx;
sc_status rc = sc_init(&ctx, SC_AES_128, SC_BACKEND_AUTO,
                       key, sizeof(key));
rc = sc_encrypt_blocks(&ctx, input, output, blocks);
\end{lstlisting}

ARM64 检测 AES、PMULL 与 SM4 能力；x86 检测 AES-NI、SSSE3、AVX2、
VAES、PCLMUL、VPCLMUL 和 GFNI。显式请求不可用 ISA 时返回
\texttt{SC\_ERR\_UNSUPPORTED}，不做静默冒充。

\section{优化技术}

\subsection{T-table：空间换指令}

\begin{table}[htbp]
\centering
\caption{表布局与数据路径}
\begin{tabular}{lll}
\toprule
算法/后端 & 表大小 & 每轮方法\\
\midrule
AES 四表 & 4~KiB & SubBytes+ShiftRows+MixColumns 合并\\
SM4 四表 & 4~KiB & 四个字节分别查表后异或\\
SM4 单表 & 1~KiB & 查一个表，再旋转 8/16/24 bit\\
SM4 重叠 & 2~KiB & 8 字节重复项，偏移 0/3/2/1 读取\\
\bottomrule
\end{tabular}
\end{table}

表越小不保证越快：1~KiB 路径增加旋转，2~KiB 路径增加非对齐字节组装；
4~KiB 路径则给 L1 数据缓存更大压力。四条路径同时存在，性能结论交给同一
基准而不是先验判断。

\subsection{Shuffle、bitslice 与 fixslice}

TWINE 的 4-bit S 盒天然适合 16-byte TBL：一条 TBL 并行完成 16 个 nibble
查找，第二次 TBL 完成轮置换。SM4 的 8-bit S 盒不能直接由单个 16 项表
表示，本项目使用 16 行固定表和掩码选择，访问轨迹与秘密无关；这是常量时间
教学路径，但单块时指令数高。GIFT 的布尔 S 盒和固定 bit permutation
完全不查秘密索引表，更接近 bitslice/fixslice 思路。

\subsection{新指令集}

\begin{itemize}
  \item \textbf{AES：}ARM AESE/AESD 为本机四路运行路径；x86 AES-NI
  与 VAES 已进入四路/八路实际后端并通过反汇编确认。
  \item \textbf{GHASH：}PMULL 做两个 64-bit 多项式的 Karatsuba 乘法，
  再以 $x^{128}=x^7+x^2+x+1$ 两次折叠约减；x86 PCLMUL/VPCLMUL
  已进入实际 GHASH 后端。
  \item \textbf{GFNI：}SM4 S 盒按
  $Y=A_1X+C_1,\ S(X)=A_2Y^{-1}+C$ 分解。标量 GFNI 模型穷举 256 个
  输入并逐项等于标准 S 盒，x86 源码用两条 affine/affine-inverse 指令。
  \item \textbf{AES 辅助 SM4：}先用固定扫描预映射到 AES 逆 S 盒值，
  再以 AESE/AESENC 完成四块 SM4 S 盒并逆转 ShiftRows；全程不以秘密值
  直接索引缓存。
  \item \textbf{专用 SM4：}ARM SM4E/SM4EKEY 与 Intel
  VSM4RNDS4/VSM4KEY4 已接入运行后端并检查实际函数；M2 Pro
  不支持 SM4 指令，运行时跳过。
\end{itemize}

\section{CTR、GCM 与 XTS}

\subsection{CTR}

CTR 对四类密码全部开放。计数器按分组宽度大端递增，每批生成八个计数器并
调用多块 API；末尾不足一块时只异或所需字节。调用前先验证全部计数器空间，
回绕时不修改输出，并返回错误码 \texttt{SC\_ERR\_COUNTER\_WRAP}。

\subsection{GCM}

GCM 只接受 128-bit 分组的 AES/SM4。96-bit IV 直接形成
$J_0=\mathrm{IV}\|0^{31}\|1$，其他 IV 先经 GHASH。AAD、密文和长度字段
均按 SP~800-38D 处理。解密先以常量时间比较标签；失败返回
\texttt{SC\_ERR\_AUTH} 并清零输出。NIST AES-GCM 与 RFC~8998
SM4-GCM 向量均纳入回归。GHASH 预计算 $H^2,H^3,H^4$，完整的四块组用
四个无数据依赖的乘法聚合。

\subsection{XTS}

XTS 使用两个独立密钥上下文，拒绝短于 16 字节的数据单元，并对非整块尾部
执行 ciphertext stealing。IEEE~1619 约定在每个字节内左移并以
\texttt{0x87} 约减；GB/T~17964-2021 约定在每个字节内右移并以
\texttt{0xe1} 约减。OpenSSL~3.6 的 SM4-XTS 默认值为 GB，本项目同时
保留 IEEE 选项。普通分组按八路预生成 tweak 并批量执行 XEX；CTS 在写回前
保存尾部，因而支持原地加解密。

\section{测试与工程验证}

\begin{table}[htbp]
\centering
\caption{验证层次}
\begin{tabularx}{\textwidth}{lX}
\toprule
层次 & 内容\\
\midrule
KAT & AES 三种密钥、SM4、GIFT-64、TWINE-80/128、NIST GCM、RFC 8998\\
内部交叉 & 参考/表/shuffle/硬件后端，固定种子随机输入，原地与非对齐语义\\
模式边界 & CTR 尾部/回绕、GCM AAD/错误标签清零、XTS CTS/最小数据单元\\
OpenSSL & x86 上 \XOpenSSLChecks 项 AES/SM4 ECB、CTR、GCM、XTS
随机差分；覆盖全部可用后端\\
工具链 & \texttt{-Wall -Wextra -Wpedantic -Werror}、ASan、UBSan\\
ISA & ARM AESE/TBL/PMULL/SM4E；\newline
x86 AES/VAES/PSHUFB/GFNI/CLMUL/VSM4\\
\bottomrule
\end{tabularx}
\end{table}

\section{实验方法}

ARM 实验平台为 \BenchCPU（\BenchPlatform），编译器为 \BenchCompiler；
x86 实验平台为 \XBenchCPU（\XBenchPlatform），编译器为
\XBenchCompiler，OpenSSL 为 \XOpenSSLVersion。
消息长度取 64~B、512~B、4~KiB、8~KiB、64~KiB、1~MiB。每组预热
3 次、正式运行 15 次；每个样本至少处理 256~KiB，记录完整纳秒值，
汇总 median、Q1、Q3、IQR、ns/byte 和十进制 GB/s。密钥扩展在计时区间外，
模式的计数器、tweak、GHASH 和标签均在计时区间内。volatile checksum
防止编译器删除计算。

x86-64 结果同时记录单调时钟和序列化 RDTSCP。文中的 TSC ticks/byte
表示 invariant-TSC 计数，不等同于随睿频变化的物理核心周期。ARM 与 x86
结果分别保存，不把一个平台的数据改名为另一个平台。

\section{结果与分析}

\begin{figure}[htbp]
\centering
\includegraphics[width=\textwidth]{results/figures/block_backends_1m.pdf}
\caption{1~MiB 分组加密吞吐率。横轴为对数尺度，数据由汇总 JSON 生成。}
\end{figure}

在 1~MiB 数据上，AES 参考、四表和硬件后端分别为
\AESRefGbps、\AESTableGbps、\AESHWGbps~GB/s。T-table 相对参考为
\AESTableGain$\times$，AESE/AESMC 为 \AESHWGain$\times$。
SM4 四表由 \SMFourRefGbps 提升至 \SMFourTableGbps~GB/s，即
\SMFourTableGain$\times$。1~KiB 单表非常接近四表，2~KiB 重叠读取在
本机反而较慢，说明“更小表”必须与额外旋转/装配成本一起评价。

\begin{figure}[htbp]
\centering
\includegraphics[width=\textwidth]{results/figures/size_scaling.pdf}
\caption{不同消息长度下 AES/SM4 后端的稳态吞吐。}
\end{figure}

硬件 AES 在 64~B 时已经领先，随后趋于稳定；本实验把函数调用和循环调度
算入短消息成本。SM4 shuffle 每轮固定扫描 16 行表，并以四块并行摊薄装配
开销；其目标是展示无秘密相关访存，而不是声称单块最高性能。
TWINE TBL 与参考接近，说明仅替换 S 盒还不足以自动获得多块并行收益。

\begin{figure}[htbp]
\centering
\includegraphics[width=\textwidth]{results/figures/mode_throughput.pdf}
\caption{512~B 与 4~KiB 的 CTR/GCM/XTS 吞吐率。}
\end{figure}

AES-CTR 能直接发挥硬件块加密能力；XTS 还要执行第二个密钥的 tweak 加密与
每块乘 $\alpha$，因此低于 CTR。GCM 已用 $H^2$--$H^4$ 消除四块组内的
串行依赖，跨组仍保留必要的 GHASH 依赖。IEEE 与 GB 两种 SM4-XTS
吞吐相近，符合两者只改变轻量 tweak 乘法的预期。

\small
\begin{longtable}{lllrrr}
\caption{1~MiB 全部实测结果（直接由 JSON 生成）}\\
\toprule
算法 & 后端 & 操作 & GB/s & ns/byte & IQR\\
\midrule
\endfirsthead
\toprule
算法 & 后端 & 操作 & GB/s & ns/byte & IQR\\
\midrule
\endhead
\OneMiBResultRows
\bottomrule
\end{longtable}
\normalsize

\subsection{x86-64 原生执行结果}

x86 状态文件记录为 \texttt{passed}。CPUID/XGETBV 检测到的 AES-NI、
SSSE3、AVX2、VAES、GFNI、PCLMUL 与 VPCLMUL 均由专用测试直接调用，
并与参考实现逐项比较。GFNI 硬件路径还穷举全部 256 个 SM4 S 盒输入；
PCLMUL 和 VPCLMUL 则分别对照独立标量 GHASH oracle。VSM4 未被 CPUID
报告，因此显式请求 \texttt{sm4-hw} 返回 \texttt{SC\_ERR\_UNSUPPORTED}。

\begin{figure}[htbp]
\centering
\includegraphics[width=\textwidth]{results/figures/x86_block_backends_1m.pdf}
\caption{x86-64 上 1~MiB 分组加密吞吐率；数据来自独立 x86 汇总。}
\end{figure}

\begin{figure}[htbp]
\centering
\includegraphics[width=\textwidth]{results/figures/x86_size_scaling.pdf}
\caption{x86-64 上 AES/SM4 后端随消息长度的吞吐变化。}
\end{figure}

\begin{figure}[htbp]
\centering
\includegraphics[width=\textwidth]{results/figures/x86_mode_throughput.pdf}
\caption{x86-64 上 CTR、GCM 与 XTS 的 512~B/4~KiB 吞吐率。}
\end{figure}

\small
\begin{longtable}{lllrrr}
\caption{x86-64 上 1~MiB 全部实测结果（直接由 JSON 生成）}\\
\toprule
算法 & 后端 & 操作 & GB/s & ns/byte & TSC ticks/byte\\
\midrule
\endfirsthead
\toprule
算法 & 后端 & 操作 & GB/s & ns/byte & TSC ticks/byte\\
\midrule
\endhead
\XOneMiBResultRows
\bottomrule
\end{longtable}
\normalsize

\section{安全、性能与平台权衡}

\begin{enumerate}
  \item \textbf{T-table：}通常减少指令，却泄露秘密相关缓存地址；只用于教学。
  \item \textbf{Shuffle/bitslice：}访问模式固定，更适合常量时间；需要足够
  多并行块才能摊薄转置和数据装配。
  \item \textbf{硬件密码指令：}AES/PMULL 兼顾性能和固定执行路径，但仍不能
  替代协议级 nonce 管理、密钥生命周期和故障防护。
  \item \textbf{运行时分派：}可移植二进制必须检查 CPUID/XGETBV 或 ARM
  feature；“编译器接受 intrinsic”不等于当前 CPU 可以执行。
  \item \textbf{x86 平台：}AES-NI/VAES、SSSE3、GFNI 与 CLMUL 已完成
  原生执行和性能采样；VSM4 仅验证编码与分派拒绝，不提供伪造吞吐。
\end{enumerate}

\section{复现方式}

\begin{lstlisting}[language=bash]
make test
make test-sanitize
make check-x86
make validate-x86
make bench
make report
\end{lstlisting}

数据链完全由文件驱动：
\begin{itemize}
  \item 原始样本：\path{results/raw/native_samples.csv}；
  \item 汇总统计：\path{results/summary/native_summary.json}；
  \item x86 原始样本：\path{results/raw/x86_samples.csv}；
  \item x86 汇总：\path{results/summary/x86_summary.json}；
  \item x86 状态：\path{results/summary/x86_status.json}；
  \item 图表脚本：\path{scripts/plot_results.py}；
  \item 数值宏脚本：\path{scripts/generate_report_data.py}。
\end{itemize}
因此每个数字都能回溯到正式样本。

\section{结论}

本项目建立了从正确性基线到多类优化和三种工作模式的完整链条。实测说明：
表优化的收益受缓存与额外指令共同约束；硬件 AES 的数量级优势最稳定；
shuffle/fixslice 必须和多块数据布局一起设计；GCM/XTS 不能只优化分组密码
而忽略 GHASH、计数器和 tweak。ARM64 与 x86-64 均已完成原生验证；
x86 结果额外证明 GFNI、VAES、PCLMUL/VPCLMUL 的实际执行，并以独立文件
保存 RDTSCP TSC ticks/byte。VSM4 仍严格限定为目标硬件不支持的静态证据。

\begin{thebibliography}{9}\sloppy
\bibitem{fips197} NIST, \emph{FIPS 197: Advanced Encryption Standard}, 2023.
\url{https://csrc.nist.gov/pubs/fips/197/final}
\bibitem{sm4} 国家标准化管理委员会，\emph{GB/T 32907-2016 SM4 分组密码算法}.
\url{https://openstd.samr.gov.cn/bzgk/std/newGbInfo?hcno=7803DE42D3BC5E80B0C3E5D8E873D56A}
\bibitem{gift} Banik et al., \emph{GIFT: A Small Present}, CHES 2017.
\url{https://www.iacr.org/archive/ches2017/10529227/10529227.pdf}
\bibitem{twine} Suzaki et al., \emph{TWINE: A Lightweight Block Cipher for Multiple Platforms}, SAC 2012.
\url{https://jpn.nec.com/rd/tg/dss/pdf/twine_SAC_full_v5.pdf}
\bibitem{modes} NIST, \emph{Block Cipher Modes}.
\url{https://csrc.nist.gov/projects/block-cipher-techniques/bcm}
\bibitem{rfc8998} P. Yang, \emph{RFC 8998: ShangMi Cipher Suites for TLS 1.3}, 2021.
\url{https://www.rfc-editor.org/rfc/rfc8998.html}
\bibitem{gfni} W. Guo, \emph{Efficient Constant-Time Implementation of SM4}, 2022/1154.
\url{https://eprint.iacr.org/2022/1154.pdf}
\bibitem{opensslxts} OpenSSL, \emph{EVP\_EncryptInit: xts\_standard}.
\url{https://docs.openssl.org/3.4/man3/EVP_EncryptInit/}
\end{thebibliography}

\end{document}
